\chapter{Elementary Potential Problems}

In quantum mechanics, we often approximate complex potentials with simpler forms. For instance, a smooth transition potential:

\begin{equation}
W(x)=\frac{V_{0}}{2}\left(1+\tanh \left(\frac{x}{d}\right)\right)
\end{equation}

Can be modeled as a step potential:

\begin{equation}
V(x)=V_{0} \theta(x)
\end{equation}

Where the Heaviside function is defined as:

\[
\theta(x)= \begin{cases}0 & \text { for } x<0  \\ 1 & \text { for } x \geq 0\end{cases}
\]

Two fundamental potential models are particularly useful:

The finite potential well:

\[
V_{-}(x)= \begin{cases}0 & \text { for } x<-a  \\ -V_{0} & \text { for }-a \leq x \leq a \\ 0 & \text { for } x>a\end{cases}
\]

This approximates potentials like:

\begin{equation}
W(x)=\frac{-V_{0}}{\cosh \left(\frac{x}{d}\right)}
\end{equation}

And the finite potential barrier:

\[
V_{+}(x)= \begin{cases}0 & \text { for } x<-a  \\ V_{0} & \text { for }-a \leq x \leq a \\ 0 & \text { for } x>a\end{cases}
\]

Which models potentials of the form:

\begin{equation}
W(x)=\frac{V_{0}}{\cosh \left(\frac{x}{d}\right)}
\end{equation}

Here's the revised LaTeX code with a different structure while maintaining all mathematical expressions:


\section{Step Potential: Solving the Eigenvalue Problem}

The step potential divides space into two regions with different energy characteristics. Starting with the Hamiltonian eigenvalue equation:

\begin{gather}
\hat{H} \Psi=E \Psi  \\
-\frac{\hbar^{2}}{2 m} \frac{d^{2} \Psi}{d x^{2}}+V_{0} \theta(x) \Psi=E \Psi
\end{gather}

We can rewrite this as a piecewise differential equation system:

\[
\begin{cases}-\frac{\hbar^{2}}{2 m} \frac{\partial^{2} \Psi}{\partial x^{2}}=E \Psi, & x<0 \text { (region I) }  \\ -\frac{\hbar^{2}}{2 m} \frac{\partial^{2} \Psi}{\partial x^{2}}+V_{0} \Psi=E \Psi, & x \geq 0 \text { (region II) }\end{cases}
\]

For the first region, we obtain a free-particle solution:

\begin{align}
& \text { Let } k=\sqrt{\frac{2 m E}{\hbar^{2}}}  \\
& \Psi_{1}=A \mathrm{e}^{i k x}+R \mathrm{e}^{-i k x}
\end{align}

The coefficient A corresponds to the incident wave, while R represents the reflected component.

For the second region, unlike classical mechanics where total reflection occurs when $E < V_0$, quantum mechanics allows for different behaviors. The general solution is:

\begin{align}
& \text { Let } q=\sqrt{\frac{2 m\left(E-V_{0}\right)}{\hbar^{2}}}  \\
& \Psi_{2}=T \mathrm{e}^{i q x}
\end{align}

No wave from the right is included as this would be physically unrealistic.

Verifying our solutions by substitution into the Schrödinger equations:

\begin{align}
k & =\sqrt{\frac{2 m E}{\hbar^{2}}}  \\
q & =\sqrt{\frac{2 m\left(E-V_{0}\right)}{\hbar^{2}}}
\end{align}

The complete wavefunction can be expressed as:

\begin{equation}
\Psi=\Psi_{1} \theta(-x)+\Psi_{2} \theta(x)
\end{equation}

To determine the coefficients, we apply boundary conditions:

\begin{enumerate}
  \item Wavefunction continuity: $\Psi_{1}(0)=\Psi_{2}(0)$
  \item Derivative continuity: $\Psi_{1}^{\prime}(0)=\Psi_{2}^{\prime}(0)$
  \item Probability current continuity: $j_{1}(0)=j_{2}(0)$
\end{enumerate}

From the first condition:

\begin{equation}
A+R=T
\end{equation}

From the second condition:

\begin{equation}
i k A-i k R=i q T
\end{equation}

Solving these equations yields:

\begin{align}
\frac{R}{A} & =\frac{k-q}{k+q} \\
\frac{T}{A} & =\frac{2 k}{k+q}
\end{align}

Substituting into our wavefunction gives:

\begin{equation}
\Psi(x)=A\left[\left(e^{i k x}+\frac{k-q}{k+q} e^{-i k x}\right) \theta(-x)+\frac{2 k}{k+q} e^{i q x} \theta(x)\right]
\end{equation}


\section{Probability Current Analysis for Step Potential}

To verify the condition $j_{1}(0)=j_{2}(0)$, we need to calculate the probability currents in both regions. Beginning with region I:

\begin{equation}
j_{1}=\frac{\hbar}{2 m i}\left(\Psi_{1}^{*} \frac{\partial \Psi_{1}}{\partial x}-\Psi_{1} \frac{\partial \Psi_{1}^{*}}{\partial x}\right)
\end{equation}

We evaluate each component individually:

\begin{align}
\Psi_{1} & =A e^{i k x}+R e^{-i k x} \\
\frac{\partial \Psi_{1}}{\partial x} & =i k A e^{i k x}-i k R e^{-i k x} \\
\Psi_{1}^{*} & =A e^{-i k x}+R e^{i k x}  \\
\frac{\partial \Psi_{1}^{*}}{\partial x} & =-i k A e^{-i k x}+i k R e^{i k x}
\end{align}

Substituting these expressions into the current formula:

\begin{align}
j_{1}= & \frac{\hbar}{2 m i}\left[\left(A e^{-i k x}+R e^{i k x}\right)\left(i k A e^{i k x}-i k R e^{-i k x}\right)\right. \\
& \left.-\left(A e^{i k x}+R e^{-i k x}\right)\left(-i k A e^{-i k x}+i k R e^{i k x}\right)\right]= \\
= & \frac{\hbar}{2 m i}\left[i k A^{2}-i k A R e^{-2 i k x}+i k R A e^{2 i k x}-i k R^{2}\right.  \\
& \left.+i k A^{2}-i k A R e^{2 i k x}+i k R A e^{-2 i k x}-i k R^{2}\right]
\end{align}

The oscillatory terms cancel out, yielding:

\begin{equation}
j_{1}=\frac{\hbar k}{m}\left(A^{2}-R^{2}\right)
\end{equation}

For region II, we compute:

\begin{equation}
j_{2}=\frac{\hbar}{2 m i}\left(\Psi_{2}^{*} \frac{\partial \Psi_{2}}{\partial x}-\Psi_{2} \frac{\partial \Psi_{2}^{*}}{\partial x}\right)
\end{equation}

With the following components:

\begin{align}
\Psi_{2} & =T e^{i q x} \\
\frac{\partial \Psi_{2}}{\partial x} & =i q T e^{i q x}  \\
\Psi_{2}^{*} & =T e^{-i q x} \\
\frac{\partial \Psi_{2}^{*}}{\partial x} & =-i q T e^{-i q x}
\end{align}

Inserting these into the current expression:

\begin{align}
j_{2} & =\frac{\hbar}{2 m i}\left[\left(T e^{-i q x}\right)\left(i q T e^{i q x}\right)-\left(T e^{i q x}\right)\left(-i q T e^{-i q x}\right)\right] \\
& =\frac{\hbar}{2 m i}\left[i q T^{2}-\left(-i q T^{2}\right)\right]  \\
& =\frac{\hbar}{2 m i}\left(2 i q T^{2}\right) \\
& =\frac{\hbar q}{m} T^{2}
\end{align}

The probability current can be divided into three distinct contributions:

\begin{align}
j_{i n} & =\frac{\hbar k}{m} A^{2} \\
j_{R} & =\frac{\hbar k}{m} R^{2}  \\
j_{T} & =\frac{\hbar q}{m} T^{2}
\end{align}

This leads to a conservation principle analogous to Kirchhoff's law:

\begin{equation}
j_{1}=j_{2} \Longrightarrow j_{i n}=j_{R}+j_{T}
\end{equation}

We can express the reflection and transmission ratios as:

\begin{align}
\frac{j_{R}}{j_{\text {in }}} & =\frac{R^{2}}{A^{2}}=\left(\frac{k-q}{k+q}\right)^{2} \\
\frac{j_{T}}{j_{\text {in }}} & =\frac{q}{k} \frac{T^{2}}{A^{2}}=\frac{q}{k}\left(\frac{2 k}{k+q}\right)^{2}=\frac{4 k q}{(k+q)^{2}}
\end{align}

For the high-energy limit where $E \gg V_{0}$:

\begin{equation}
q=\sqrt{\frac{2 m\left(E-V_{0}\right)}{\hbar^{2}}} \simeq \sqrt{\frac{2 m E}{\hbar^{2}}}=k
\end{equation}


\section{Analysis of Energy Regimes for Step Potential}

For the high-energy case ($E \gg V_{0}$), the transmission and reflection ratios become:

\begin{align}
& \frac{j_{R}}{j_{i n}}=\frac{R^{2}}{A^{2}}=\left(\frac{k-q}{k+q}\right)^{2} \simeq 0  \\
& \frac{j_{T}}{j_{i n}}=\frac{4 k q}{(k+q)^{2}} \simeq \frac{4 k^{2}}{4 k^{2}}=1
\end{align}

This physical result aligns with our intuition: particles with energy significantly exceeding the barrier height pass through virtually unimpeded. The wavefunction simplifies to:

\begin{equation}
\Psi \approx A \mathrm{e}^{i k x}
\end{equation}

When examining the threshold case where $E \simeq V_{0}$:

\begin{equation}
q=\sqrt{\frac{2 m\left(E-V_{0}\right)}{\hbar^{2}}} \simeq 0
\end{equation}

The transmission and reflection ratios become:

\begin{align}
\frac{j_{R}}{j_{\text {in }}} & =\frac{R^{2}}{A^{2}}=\left(\frac{k-q}{k+q}\right)^{2} \simeq \frac{k^{2}}{k^{2}}=1  \\
\frac{j_{T}}{j_{\text {in }}} & =\frac{4 k q}{(k+q)^{2}} \simeq 0
\end{align}

In this scenario, the wave experiences complete reflection with no transmission. The amplitude coefficients satisfy:

\begin{align}
& \frac{R}{A}=\frac{k-q}{k+q} \simeq 1 \Longrightarrow R=A  \\
& \frac{T}{A}=\frac{2 k}{k+q} \simeq 2 \Longrightarrow T=2 A
\end{align}

The regional wavefunctions become:

\begin{equation}
\Psi_{1}=A \mathrm{e}^{i k x}+A \mathrm{e}^{-i k x}=2 A \cos (k x)
\end{equation}

And:

\begin{equation}
\Psi_{2}=T \mathrm{e}^{i q x} \simeq T=2 A
\end{equation}

For the classically forbidden region where $E<V_{0}$, we redefine $q$ to maintain a real parameter:

\begin{equation}
q=\sqrt{\frac{2 m\left(V_{0}-E\right)}{\hbar}}
\end{equation}

The wavefunctions take the form:

\begin{align}
& \Psi_{1}=A \mathrm{e}^{i k x}+R \mathrm{e}^{-i k x} \\
& \Psi_{2}=T \mathrm{e}^{-q x}
\end{align}

The boundary conditions yield:

\[
\left\{\begin{array}{l}
A+R=T  \\
i k A-i k R=-q T
\end{array}\right.
\]


\section{Derivation of Coefficients for Sub-barrier Case}

Solving the boundary condition equations for the case where $E<V_0$:

\begin{align}
& \left\{\begin{array}{l}
A+R=T \\
k R-k A=-i q A+-i q R
\end{array}\right. \\
& \left\{\begin{array}{l}
A+R=T \\
R(k+i q)=A(k-i q)
\end{array}\right. \\
& \left\{\begin{array}{l}
A+R=T \\
\frac{R}{A}=\frac{k-i q}{k+i q}
\end{array}\right.  \\
& \left\{\begin{array}{l}
1+\frac{R}{A}=\frac{T}{A} \\
\frac{R}{A}=\frac{k-i q}{k+i q}
\end{array}\right. \\
& \left\{\begin{array}{l}
\frac{k-i q+k+i q}{k+i q}=\frac{T}{A} \\
\frac{R}{A}=\frac{k-i q}{k+i q}
\end{array}\right.
\end{align}

This yields the amplitude ratios:

\begin{align}
\frac{R}{A} & =\frac{k-i q}{k+i q} \\
\frac{T}{A} & =\frac{2 k}{k+i q}
\end{align}

The complete wavefunction becomes:

\begin{equation}
\Psi(x)=A\left[\left(e^{i k x}+\frac{k-i q}{k+i q} e^{-i k x}\right) \theta(-x)+\frac{2 k}{k+i q} e^{-q x} \theta(x)\right]
\end{equation}

For region II, we analyze the probability current:

\begin{equation}
j_{2}=\frac{\hbar}{2 m i}\left(\Psi_{2}^{*} \frac{\partial \Psi_{2}}{\partial x}-\Psi_{2} \frac{\partial \Psi_{2}^{*}}{\partial x}\right)
\end{equation}

Evaluating each term:

\begin{align}
\Psi_{2} & =T e^{-q x} \\
\frac{\partial \Psi_{2}}{\partial x} & =-q T e^{-q x} \\
\Psi_{2}^{*} & =T^{*} e^{-q x}  \\
\frac{\partial \Psi_{2}^{*}}{\partial x} & =-q T^{*} e^{-q x}
\end{align}

Substituting into the current expression:

\begin{align}
j_{2} & =\frac{\hbar}{2 m i}\left[\left(T e^{-q x}\right)\left(-q T^{*} e^{-q x}\right)-\left(T e^{-q x}\right)\left(-q T^{*} e^{-q x}\right)\right] \\
& =\frac{\hbar}{2 m i}\left[-q|T|^{2} e^{-2 q x}+q|T|^{2} e^{-2 q x}\right]=0
\end{align}

Since $j_2 = 0$, the conservation law requires $j_1 = 0$ as well. For region I:

\begin{equation}
j_{1}=\frac{\hbar}{2 m i}\left(\Psi_{1}^{*} \frac{\partial \Psi_{1}}{\partial x}-\Psi_{1} \frac{\partial \Psi_{1}^{*}}{\partial x}\right)
\end{equation}

With components:

\begin{align}
\Psi_{1} & =A e^{i k x}+R e^{-i k x} \\
\frac{\partial \Psi_{1}}{\partial x} & =i k A e^{i k x}-i k R e^{-i k x}  \\
\Psi_{1}^{*} & =A^{*} e^{-i k x}+R^{*} e^{i k x} \\
\frac{\partial \Psi_{1}^{*}}{\partial x} & =-i k A^{*} e^{-i k x}+i k R^{*} e^{i k x}
\end{align}

The current calculation yields:

\begin{align}
j_{1}= & \frac{\hbar}{2 m i}\left[\left(A^{*} e^{-i k x}+R^{*} e^{i k x}\right)\left(i k A e^{i k x}-i k R e^{-i k x}\right)\right. \\
& \left.-\left(A e^{i k x}+R e^{-i k x}\right)\left(-i k A^{*} e^{-i k x}+i k R^{*} e^{i k x}\right)\right] \\
= & \frac{\hbar}{2 m i}\left[i k|A|^{2}-i k A^{*} R e^{-2 i k x}+i k R^{*} A e^{2 i k x}-i k|R|^{2}\right. \\
& \left.+i k|A|^{2}-i k A R^{*} e^{2 i k x}+i k R A^{*} e^{-2 i k x}-i k|R|^{2}\right] \\
= & \frac{\hbar}{2 m i}\left[2 i k|A|^{2}-2 i k|R|^{2}+i k\left(R^{*} A-A R^{*}\right) e^{2 i k x}+i k\left(A^{*} R-R A^{*}\right) e^{-2 i k x}\right] \\
= & \frac{\hbar}{2 m i}\left[2 i k\left(|A|^{2}-|R|^{2}\right)+i k\left(\left(R^{*} A-A R^{*}\right) e^{2 i k x}+\left(A^{*} R-R A^{*}\right) e^{-2 i k x}\right)\right]= \\
= & \frac{\hbar k}{m}\left(|A|^{2}-|R|^{2}\right)
\end{align}


\section{Current Conservation and Physical Interpretation}

Taking $A$ as real, the current conservation condition gives:

\begin{equation}
j_{1}=\frac{\hbar k}{m}\left(A^{2}-|R|^{2}\right)=0 \Longrightarrow|A|=|R|
\end{equation}

This equality between incident and reflected amplitudes indicates complete reflection for the sub-barrier case.

\section{Finite Potential Barrier Analysis}

For a rectangular potential barrier, we divide space into three distinct regions:

\[
V= \begin{cases}0 & \text { for } x<-a  \\ V_{0} & \text { for }-a \leq x \leq a \\ 0 & \text { for } x>a\end{cases}
\]

The Schrödinger equation takes different forms in each region:

\[
\begin{cases}-\frac{\hbar^{2}}{2 m} \frac{\partial^{2} \Psi_{1}}{\partial x^{2}}=E \Psi_{1}, & x<-a \quad \text { (region I) }  \\ -\frac{\hbar^{2}}{2 m} \frac{\partial^{2} \Psi_{2}}{\partial x^{2}}+V_{0} \Psi_{2}=E \Psi_{2}, & -a \leq x \leq a \quad \text { (region II) } \\ -\frac{\hbar^{2}}{2 m} \frac{\partial^{2} \Psi_{3}}{\partial x^{2}}=E \Psi_{3}, & x>a \quad \text { (region III) }\end{cases}
\]

The general solutions in each region are:

\[
\begin{array}{ll}
\Psi_{1}(x)=A e^{i k x}+R e^{-i k x} & \text { for } x<-a \\
\Psi_{2}(x)=C e^{i q x}+D e^{-i q x} & \text { for }-a \leq x \leq a  \\
\Psi_{3}(x)=T e^{i k x} \quad \text { for } x>a &
\end{array}
\]

The wavevectors are defined as:

\begin{align}
k & =\sqrt{\frac{2 m E}{\hbar^{2}}} \\
q & =\sqrt{\frac{2 m\left(E-V_{0}\right)}{\hbar^{2}}}
\end{align}

At the boundaries between regions, we apply matching conditions:

\begin{align}
& \Psi_{1}(-a)=\Psi_{2}(-a) \\
& \Psi_{1}^{\prime}(-a)=\Psi_{2}^{\prime}(-a) \\
& \Psi_{2}(a)=\Psi_{3}(a)  \\
& \Psi_{2}^{\prime}(a)=\Psi_{3}^{\prime}(a)
\end{align}

These conditions yield the following equations:

\begin{equation}
A e^{-i k a}+R e^{i k a}=C e^{-i q a}+D e^{i q a}
\end{equation}

\begin{equation}
-i k\left(A e^{-i k a}-R e^{i k a}\right)=-i q\left(C e^{-i q a}-D e^{i q a}\right)
\end{equation}

\begin{equation}
C e^{i q a}+D e^{-i q a}=T e^{i k a}
\end{equation}

\begin{equation}
i q\left(C e^{i q a}-D e^{-i q a}\right)=i k T e^{i k a}
\end{equation}

Solving this system of equations, we obtain the reflection and transmission amplitudes:

\begin{equation}
\frac{R}{A}=\mathrm{e}^{-2 i k a} \frac{i\left(q^{2}-k^{2}\right) \sin (2 q a)}{2 k q \cos (2 q a)-i\left(k^{2}+q^{2}\right) \sin (2 q a)}
\end{equation}

\begin{equation}
\frac{T}{A}=\frac{2 k q e^{-2 i k a}}{2 k q \cos (2 q a)-i\left(k^{2}+q^{2}\right) \sin (2 q a)}
\end{equation}


\section{Probability Current Analysis for Barrier Problem}

For the finite barrier, we can express the probability currents in each region. For region I:

\begin{equation}
j_{1}=\frac{\hbar k}{m}\left(A^{2}-|R|^{2}\right)
\end{equation}

While in region III:

\begin{equation}
j_{3}=\frac{\hbar k}{m}|T|^{2}
\end{equation}

Current conservation requires $j_1 = j_2 = j_3$, leading to:

\begin{equation}
A^{2}=|R|^{2}+|T|^{2}
\end{equation}

The probability current components can be identified as:

\begin{align}
j_{i n} & =\frac{\hbar k}{m} A^{2} \\
j_{R} & =\frac{\hbar k}{m}|R|^{2}  \\
j_{T} & =\frac{\hbar k}{m}|T|^{2}
\end{align}

The reflection and transmission coefficients become:

\begin{align}
\frac{j_{R}}{j_{i n}} & =\frac{|R|^{2}}{A^{2}}=\frac{\left(q^{2}-k^{2}\right) \sin ^{2}(2 q a)}{4 k^{2} q^{2}+\left(k^{2}-q^{2}\right) \sin ^{2}(2 q a)} \\
\frac{j_{T}}{j_{i n}} & =\frac{|T|^{2}}{A^{2}}=\frac{4 k^{2} q^{2}}{4 k^{2} q^{2}+\left(k^{2}-q^{2}\right) \sin ^{2}(2 q a)}
\end{align}

\section{Limiting Cases for Barrier Transmission}

In the high-energy limit where $E \gg V_0$, we find $q \approx k$, resulting in:

\begin{align}
& \frac{R}{A} \rightarrow 0 \\
& \frac{T}{A} \rightarrow 1
\end{align}

This corresponds to complete transmission through the barrier.

As $E$ approaches $V_0$ from above ($E \rightarrow V_0^+$):

\begin{align}
& \frac{R}{A} \rightarrow \frac{-i k a}{1-i k a} \Longrightarrow \frac{j_{R}}{j_{i n}}=\frac{k^{2} a^{2}}{1+k^{2} a^{2}}  \\
& \frac{T}{A} \rightarrow \frac{1}{1-i k a} \Longrightarrow \frac{j_{R}}{j_{i n}}=\frac{1}{1+k^{2} a^{2}}
\end{align}

The behavior depends critically on the dimensionless parameter $ka$, which relates to the ratio of barrier width to wavelength:

\begin{equation}
k a=\frac{2 \pi a}{\lambda}
\end{equation}

This gives rise to two distinct regimes:

\[
\begin{array}{ll}
k a \gg 1 & \text { thick barrier } \\
k a \ll 1 & \text { thin barrier }
\end{array}
\]

For a thick barrier ($ka \gg 1$), the reflection dominates:

\begin{align}
\frac{R}{A} & =\frac{-i k a}{1-i k a} \rightarrow 1 \\
\frac{T}{A} & =\frac{1}{1-i k a} \rightarrow 0
\end{align}

For a thin barrier ($ka \ll 1$), transmission prevails:

\begin{align}
\frac{R}{A} & =\frac{-i k a}{1-i k a} \rightarrow 0 \\
\frac{T}{A} & =\frac{1}{1-i k a} \rightarrow 1
\end{align}

A particularly interesting case occurs when $q=\frac{\pi n}{a}$, where the reflection coefficient vanishes completely:

\begin{align}
\frac{R}{A} & =0 \\
\frac{T}{A} & =1
\end{align}


\section{Quantum Tunneling Analysis}

When the particle energy is below the potential barrier height ($E < V_0$), classical physics predicts total reflection. However, quantum mechanics allows for the tunneling effect. To analyze this phenomenon, we redefine the parameter $q$ as:

\begin{equation}
q \rightarrow i q
\end{equation}

This substitution transforms our previous formulas through the following identities:

\begin{equation}
\sin (i x)=i \sinh (x)
\end{equation}

\begin{equation}
\cos (i x)=\cosh (x)
\end{equation}

The wavefunctions in the three regions become:

\begin{align}
& \Psi_{1}(x)=A e^{i k x}+R e^{-i k x} \quad \text { for } x<-a \\
& \Psi_{2}(x)=C e^{q x}+D e^{-q x} \quad \text { for }-a \leq x \leq a  \\
& \Psi_{3}(x)=T e^{i k x} \quad \text { for } x>a
\end{align}

Applying boundary conditions yields the modified reflection and transmission amplitudes:

\begin{align}
\frac{R}{A} & =\mathrm{e}^{-2 i k a} \frac{\left(q^{2}-k^{2}\right) \sinh (2 q a)}{2 i k q \cosh (2 q a)+\left(k^{2}-q^{2}\right) \sinh (2 q a)} \\
\frac{T}{A} & =\frac{2 i k q e^{-2 i k a}}{2 i k q \cosh (2 q a)+\left(k^{2}-q^{2}\right) \sinh (2 q a)}
\end{align}

The corresponding probability current ratios are:

\begin{align}
\frac{j_{R}}{j_{i n}} & =\frac{\left(q^{2}+k^{2}\right)^{2} \sinh ^{2}(2 q a)}{4 k^{2} q^{2}+\left(q^{2}+k^{2}\right)^{2} \sinh ^{2}(2 q a)}  \\
\frac{j_{T}}{j_{i n}} & =\frac{4 k^{2} q^{2}}{4 k^{2} q^{2}+\left(q^{2}+k^{2}\right)^{2} \sinh ^{2}(2 q a)}
\end{align}

In the limit where the barrier height greatly exceeds the particle energy ($V_0 \gg E$), we approach total reflection:

\begin{align}
& \frac{R}{A} \simeq 1  \\
& \frac{T}{A} \simeq 0
\end{align}
